\section{Reduction to pressure increments}\label{sec:reduction}

The comparison estimates and finite pressure expansions will yield one
increment bound. We state it first and prove here that it implies both
moment convergence and the characterization in Proposition~\ref{def:law}.
The proof of the increment bound is completed in
Sections~\ref{sec:expansion} and~\ref{sec:parameters}.

Take the coordinates of \(\gamma\) independently and uniformly in
\([1/3,2/3]\), as in \eqref{eq:P3}, and let
\(\ell\in[1,2^{1/6}]\). A direct strength sweep varies
\(\log\lambda\) in \eqref{eq:P1}, initially with
\((\chi_{\rm lin},\chi_e)=(0,1)\). At final strength, two linear
multiplier sweeps pass successively to \((1,1)\) and \((1,0)\).
A dot denotes differentiation in the parameter of the current segment,
always at fixed integer site count \(m\). The size sweep uses the
smoothed path at final strength with multipliers \((1,0)\), and its dot
is \(\partial_\theta\). For a fixed order \(r\), let \(c_0,\ldots,c_r\)
be the interpolation-differentiation stencil characterized by
\[
 \sum_{j=0}^r c_j P(j)=P'(0)
 \qquad\text{for every polynomial \(P\) of degree at most \(r\)}.
\]

\begin{proposition}[Stopped pressure increments]\label{prop:increment}
There is a choice of the exponents in \eqref{eq:P1} and a constant
\(c>0\) with the following property. For every sufficiently small fixed
\(\varepsilon>0\), every fixed \(q>0\), and
\(s=q x^{1+\varepsilon}\), the following bounds hold for every
sufficiently large fixed stencil order:
\begin{enumerate}
\item On each direct segment,
\[
 \int \mathbf E_\gamma\!\left[
  {\bf1}_{\{\mathcal K_m(s)\le3x^\varepsilon\}}
       |\dot\psi_m(s)|\right]\,|d\vartheta|
 \lesssim x^{5+c},
\]
uniformly for \(N/2\le m\le3N\), where \(\vartheta\) is that
segment's parameter.
\item On the size segment the corresponding bound holds for
\begin{equation}
 \left|\dot\psi_m(s)+\sum_{j=0}^r
                   c_j(m+j)\psi_{m+j}(s)\right|.\tag{I}\label{eq:I}
\end{equation}
It also holds when \(m\) is a measurable function of \(\theta\) with
values in this integer range, or when the observables are averaged
with nonnegative probability weights on these integers depending
measurably on \(\theta\).
\end{enumerate}
All integrations use absolute path-parameter length. The saving \(c\)
can be chosen independently of sufficiently small \(\varepsilon\).
Constants may depend on the fixed test coefficient \(q\), stencil order,
and the fixed temperature.
\end{proposition}

\subsection{Exponential augmentation and the stopped change of measure}

To transfer moment estimates to the tilted law, we need fixed
\(L^p\) bounds for its density, which involve Laplace transforms at
larger fixed multiples of \(s\). Adding independent exponential
noise produces a log-concave law whose centered positive Laplace
transform controls its variance and both tails.

At fixed site count, path point, and parameters \(\gamma\), the
transform \(\mathcal K_m(s)\) is deterministic because its defining
disorder expectation has already been taken. Thus the stop
\(\mathcal K_m(s)\le3x^\varepsilon\) restricts path data. All density
norms below use the untilted initial Gaussian law.

\begin{lemma}[Augmentation along the physical paths]\label{red:lem-augmentation}
Let \(G_m\) be the continuation in \eqref{eq:A4}, and let
\(\mathcal E_{x/2}\) be an independent exponential random variable of
rate \(x/2\). Along every physical path used above,
\(G_m+\mathcal E_{x/2}\) has a log-concave density. On
\(\mathcal K_m(s)\le3x^\varepsilon\), the Radon--Nikodym density
\[
 R_m=\frac{e^{sG_m}}{\mathbf E e^{sG_m}}
\]
has bounded \(L^p\) norm for every fixed finite \(p\), uniformly in
the path and the allowed site counts for all sufficiently large \(N\).
\end{lemma}

\begin{proof}
First take a finite hierarchy and include all future Gaussian drivers
under the normalized positive law. Write the positive revelation masses
as \(v_1<\cdots<v_L\le a\), and put \(F_0=G_m\). Let \(F_i\)
be the continuation after the first \(i\) positive-mass revelations,
as a function of that Gaussian prefix. In particular \(F_L\) is the
terminal log partition function divided by \(a\). The transition at
\(v_i\) contributes \(\exp\{v_i(F_i-F_{i-1})\}\), and the terminal
Gibbs weight of a fixed spin \(\sigma\) contributes
\(\exp\{aH_\sigma-aF_L\}\). Multiplying these factors gives the
Gaussian driver density times
\[
 \exp\!\left\{
 aH_\sigma-v_1G_m
 -\sum_{i=1}^{L-1}(v_{i+1}-v_i)F_i-(a-v_L)F_L
 \right\}.
\]
Here \(H_\sigma\) is linear in the drivers and every continuation is
convex in the corresponding Gaussian prefix, hence also in the full
driver vector when extended independently of its later coordinates.
This is the telescoping
calculation in the proof of the single-spin pinned moment estimate in
\cite[Lemma~2.3]{PE}.
All physical clocks are constant between their initial revelation and
mass \(x\). Thus \(v_1\ge x\), or, if empty levels have been inserted,
the identical initial-continuation terms combine to a coefficient at
least \(x\). Rescaling mass length does not change this coefficient:
the product of a mass and its continuation is
\((v/a)(aF)=vF\). In particular the smoothed length-four path has the
same lower bound.

Before the pin, the future normalized transition kernels and terminal
Gibbs probabilities integrate to one. The initial driver therefore has
its original Gaussian law. The spin gauge transformations preserve
\(G_m\) and act transitively on terminal spins, so the law of \(G_m\)
is also unchanged by the pin. This uses the full joint positive law;
no inequality under an arbitrary conditioning is required.

For \(z=G_m+\mathcal E_{x/2}\), the independent exponential contributes
\(\exp\{(x/2)(G_m-z)\}\) on the epigraph \(G_m\le z\).
At least \(x/2\) of the negative convex-continuation coefficient
remains. The epigraph is convex, and the joint density in the drivers
and \(z\) is log-concave. Marginalization therefore gives a log-concave
density in \(z\), by the log-concave marginal theorem
\cite{Prekopa1973}. Fixed-size path approximation passes the claim to
the paths in use: the laws converge weakly, their log-concavity
inequalities pass on intervals, and the exponential convolution gives
a density and zero mass to interval endpoints.

Let
\[
 Z=G_m+\mathcal E_{x/2}
       -\mathbf E(G_m+\mathcal E_{x/2}),\qquad d_G^2=\mathbf E Z^2.
\]
The one-dimensional log-concave estimate of
\cite[Lemma~7.7]{PE}
gives two-sided exponential tails on scale \(d_G\), and constants
\(c_1>0\) such that \(\mathbf P(Z\ge c_1d_G)\ge c_1\).
On the stopped set,
\[
 \log\mathbf E e^{sZ}
 =\mathcal K_m(s)-\log(1-2s/x)-2s/x
 \le3x^\varepsilon+O(x^{2\varepsilon}).
\]
The function \(e^t-1-t\) is nonnegative, and on the displayed
right-hand event it is at least \((s c_1d_G)^2/2\). Since
\(\mathbf EZ=0\), it follows that
\[
 c_1(s c_1d_G)^2/2
 \le\mathbf E(e^{sZ}-1-sZ)
 =e^{\log\mathbf E e^{sZ}}-1
 \lesssim x^\varepsilon.
\]
Consequently \(s d_G\lesssim x^{\varepsilon/2}\). The two-sided
exponential tails now bound \(\mathbf E e^{p sZ}\) for every fixed
\(p\). Conditional Jensen over the independent centered exponential
transfers this bound to \(G_m-\mathbf E G_m\). Since
\(\mathcal K_m(s)\ge0\), this proves the asserted bound for \(R_m\).
\end{proof}

The active-component law in \eqref{eq:A4} may therefore be estimated by
H\"older from its untilted marginal, with an arbitrarily small power
loss when the required moment orders are fixed sufficiently large.
For a product involving several active components, apply these marginal
bounds and H\"older to the finite collection of density factors. Their
roots remain independent; no common untilted root is introduced.
Outside the stopped set no derivative estimate of this accuracy will
be used.

\subsection{Mollification in the number of sites}

The size path changes both the clocks and the number of sites.
At fixed integer \(m\), differentiation of
\(\mathcal K_m(s)=sm\psi_m(s)\) accounts only for the clock change.
We average over nearby integers to obtain a differentiable function
of a real size center. Differentiating the averaging weights then
produces the finite-difference term in \eqref{eq:I}; a neighboring-site
bound will let us remove the averaging at the endpoints.

Let \(\rho\) be a fixed nonnegative smooth compactly supported
probability density, and set
\[
 D=x^{-1/2},\qquad n=Ne^\theta,\qquad
 w_m(n)=D^{-1}\rho\bigl((m-n)/D\bigr).
\]
The real center \(n\) is frozen on direct segments. Smooth
sum--integral comparison gives, for every fixed \(k\),
\[
 \sum_m w_m(n)=1+O_k(D^{-k}).
\]
For example, repeated integration by parts in the Fourier series of
the periodization proves the estimate uniformly in \(n\).

\begin{lemma}[Stencil and neighboring-site errors]\label{red:lem-site-smoothing}
For any sequence \(Z_m\), supported or restricted to the relevant
neighborhood of the center,
\[
 \frac{d}{dn}\sum_m w_m(n)Z_m
 =\sum_m w_m(n)\sum_{j=0}^r c_j Z_{m+j}
       +O_r\!\left(D^{-r-1}\sup_m|Z_m|\right).
\]
At the same clock data, uniformly along the paths under consideration,
\begin{equation}\label{red:eq-neighbor}
 |\mathcal K_{m+1}(s)-\mathcal K_m(s)|\le Cs(1+M).
\end{equation}
\end{lemma}

\begin{proof}
Move the stencil from \(Z\) to the weights. The coefficient of
\(Z_m\) becomes \(\sum_jc_jw_{m-j}\). Taylor expansion of the shifted
weights, together with the defining polynomial identities of the
stencil, leaves the derivative in \(n\). Its remainder has size
\(O_r(D^{-r-2})\) for each weight and is summed over \(O(D)\) sites.

For \eqref{red:eq-neighbor}, append a zero-input independent spin and
interpolate to the completed \((m+1)\)-site model with the adjusted
matrix normalization. The change in total cumulative covariance is
\(O(1+M)\), and the affine interpolation is positive semidefinite.
Apply pressure differentiation to both tests in \eqref{eq:A4}.
The spin prior is a probability measure, so the unused spin contributes
no additional constant. This gives the displayed bound.
\end{proof}

\begin{proposition}[Dyadic transform comparison]\label{red:prop-dyadic}
Assume Proposition~\ref{prop:increment}. There is \(c'>0\), independent
of sufficiently small fixed \(\varepsilon\), such that for every fixed
positive test coefficient \(q\), every integer \(N'\in[N,2N]\), and
\(s=q x^{1+\varepsilon}\),
\[
 \left|\log\mathbb E e^{(s/x)Y_{N'}}
       -\log\mathbb E e^{(s/x)Y_N}\right|
 \lesssim x^{\varepsilon+c'},\qquad
 Y_m=m^{-1/6}(F_m^c-\mathbb EF_m^c).
\]
The constants are uniform in the integer \(N'\).
\end{proposition}

\begin{proof}
Use a direct sweep at center \(N\), the switch \eqref{eq:P11}, the size
sweep to \(\theta=\log(N'/N)\), the reverse switch, and the reverse
direct sweep at center \(N'\). At both ends the regularization strength
is \(x^{12}\); the clocks differ from the bare data
\((T/\ell^2,0)\) by \(O_\circ(x^{12})\). The self corrections cancel
in the centered transform.

Apply Fubini to the finite sum of the nonnegative stopped derivative
integrals, with the site weights above. Their total expectation in
\(\gamma\) is \(O(x^{5+c})\), so one choice of \(\gamma\) has this
bound up to a fixed constant. The weights can be normalized by their
sum, with an error smaller than any fixed power by increasing the
sum--integral order. The weighted form of \eqref{eq:I} expressly
permits their dependence on \(\theta\). The choice is made for this
one dyad comparison and one test \(s\); different tests may use
different paths because the bare endpoint laws do not depend on
\(\gamma\).

Write
\[
 \overline{\mathcal K}(\theta)=\sum_mw_m(n)\mathcal K_m(s).
\]
By \eqref{red:eq-neighbor}, its active site values have oscillation
\(O(sD(1+M))=x^{\varepsilon+1/2-o(1)}\). Since
\(\mathcal K_m\ge0\) and \(\sum_mw_m=1+O_k(D^{-k})\), whenever
\(\overline{\mathcal K}\le x^\varepsilon\), all contributing site
counts, and the boundedly many neighboring stencil counts, lie in
their individual stops \(\mathcal K_m\le3x^\varepsilon\).

On a size segment use \(Z_m=m\psi_m(s)\) in
Lemma~\ref{red:lem-site-smoothing}. The derivative is
\begin{align*}
 \frac{d}{d\theta}\overline{\mathcal K}
 ={}&sn\sum_mw_m
  \left(\dot\psi_m+\sum_{j=0}^r c_j(m+j)\psi_{m+j}\right)\\
 &+s\sum_mw_m(m-n)\dot\psi_m
   +O_\circ\bigl(sN^2D^{-r-1}\bigr).
\end{align*}
Indeed \(Z_m=O_\circ(N)\), by the pressure bounds in
\eqref{eq:A1}. The dotted-clock \(L^1\) bounds and
\eqref{eq:A1} make the middle term \(O_\circ(sD)\), after integration
as well. The two auxiliary errors have sizes
\[
 sD=x^{\varepsilon+1/2},\qquad
 sN^2D^{-r-1}=x^{\varepsilon-11+(r+1)/2}
\]
up to subpower losses. Choose the fixed stencil order large enough to
leave a positive saving. Multiplying the increment bounds by
\(sn=O(x^{-5+\varepsilon})\), and including the switches
\eqref{eq:P11}, gives total drift
\(O(x^{\varepsilon+c'})\). The bare endpoint approximation costs
\(O_\circ(sN x^{12})=O_\circ(x^{7+\varepsilon})\).
The savings can be decreased independently of \(\varepsilon\).

At the bare initial endpoint the fixed-temperature exponent theorem
of \cite[Theorem~1.1]{PE}, together with the augmented
log-concave estimate, gives
\(\mathcal K_m(s)=o(x^\varepsilon)\) uniformly in the active window.
Indeed the variance scale is \(x^{-2-o(1)}\) before multiplication
by \(x\), and the small-positive-parameter bound gives
\(O(x^{2\varepsilon-o(1)})\). Continuity on each segment, and the
small switch errors, exclude a first violation of
\(\overline{\mathcal K}\le x^\varepsilon\). Only stopped derivative
estimates have entered this argument.

Remove the endpoint weights using \eqref{red:eq-neighbor}; the cost
is again \(O(sD(1+M))\). At the final bare endpoint the terminal mass
is \(\ell=(N'/N)^{1/6}\), while the covariance is \(T/\ell^2\).
The terminal continuation is therefore distributed as
\(F_{N'}^c/\ell\), with the same physical inverse temperature.
Since \(x/\ell=N'^{-1/6}\), its \(x\)-scaled centered law is exactly
\(Y_{N'}\). The initial law is \(Y_N\). This proves the claim,
including its uniformity throughout the dyad.
\end{proof}

\subsection{Moments, positive variance, and full-sequence convergence}

\begin{proposition}[Consequences of the increment bound]\label{red:prop-moments}
Assume Proposition~\ref{prop:increment}. For every fixed integer
\(k\ge1\), the moments \(\mathbb EY_n^k\) converge at a positive
power rate. Their limiting second moment is finite and strictly
positive, and the centered variables \(Y_n\) have uniformly bounded
exponential moments on some constant neighborhood of the origin.
They converge weakly to a moment-determinate law. Removing the
completion and standardizing by the exact finite-size variance gives
the convergence asserted in Theorem~\ref{thm:main}.
\end{proposition}

\begin{proof}
Fix \(k\). The dyadic comparison gives values of the log Laplace
transform at positive arguments. We recover its derivatives at zero
by polynomial interpolation, choosing \(\varepsilon\) after \(k\)
so that the interpolation loss is smaller than the comparison saving.

At bare endpoints, the exponent theorem, the augmented tail bound,
and conditional Jensen give a scale \(L_n=x^{-o(1)}\), uniform in
the dyad, with \(\mathbb E e^{|Y_n|/L_n}\le C\). Hence, for a
sufficiently small constant \(a>0\),
\(M_n(z)=\mathbb E e^{zY_n}\) satisfies
\(|M_n(z)-1|\le1/2\) on
\(|z|\le r_n^{\mathrm{an}}:=a/L_n\), where its logarithm is analytic
and bounded by an absolute constant.
Cauchy's estimate on this disk gives
\[
 \left|\log M_n(z)-
   \sum_{j=0}^{k}\frac{(\log M_n)^{(j)}(0)}{j!}z^j\right|
 \le C_k(L_n|z|)^{k+1},
 \qquad |z|\le r_n^{\mathrm{an}}/2,
\]
which is \(x^{(k+1)\varepsilon-o(1)}\) at each fixed multiple of \(x^\varepsilon\) and uses
only the preliminary subpower scale.

Interpolate at the \(k+1\) nonnegative nodes
\(0,x^\varepsilon,\ldots,kx^\varepsilon\).
Proposition~\ref{red:prop-dyadic} bounds each nonzero-node difference;
at zero the difference is zero. Division by the interpolation scale
loses at most \(x^{-k\varepsilon}\). Thus every cumulant through
order \(k\) has dyadic difference bounded by
\[
 O\bigl(x^{c'-(k-1)\varepsilon}
              +x^{\varepsilon-o(1)}\bigr).
\]
Take \(\varepsilon\) sufficiently small after fixing \(k\), for
example with \(k\varepsilon<c'/2\). Both terms save a fixed positive
power. Summing a geometric series over adjacent dyads proves
convergence at a power rate on powers of two; the comparison with
every integer in the next dyad proves the same conclusion for all
integers. Moments are polynomials in finitely many cumulants, so they
also have power rates. Enlarge the constants over the finitely many
remaining sizes to obtain the bounds for every \(n\ge2\).

Let \(v_n=\operatorname{Var}(Y_n)\), and let its finite limit be
\(v\). If \(v=0\), its power rate would give
\(v_n=O(n^{-a})\) for some \(a>0\). The lower bound in the PE
exponent theorem gives \(v_n\ge n^{-2\eta}\) eventually for every
fixed \(\eta>0\), after a harmless adjustment of the exponent
notation. Choosing \(2\eta<a\) is a contradiction. Hence \(v>0\).

The bounded variances improve the tail scale to a uniform constant.
To see the normalization explicitly, put \(r_n=n^{-1/6}\) at a bare
system of size \(n\), and let \(E_n\) be independent exponential
noise of rate \(r_n/2\). The bare version of the pinning calculation
above gives a log-concave density for \(F_n^c+E_n\). At a bare path
the negative continuation coefficient in that calculation is the
terminal mass one, which exceeds \(r_n/2\). Thus
\[
 \widetilde Y_n=Y_n+r_n(E_n-\mathbb E E_n)
 \quad\hbox{has variance}\quad v_n+4.
\]
The log-concave tail estimate applies with a common constant scale.
Conditional Jensen, using the convex function
\(z\mapsto e^{t|z|}\), removes the centered independent noise.
Consequently, for some \(t_0>0\),
\[
 \sup_{n\ge2}\mathbb E e^{t_0|Y_n|}<\infty.
\]
Finitely many small sizes can be included by decreasing \(t_0\).
These bounds give tightness and uniform integrability of every fixed
power. They also imply moment determinacy, for example through the
resulting factorial moment bounds. Every weak subsequential limit
therefore has the same moments and is the same probability law.
Thus weak convergence holds along the full sequence.

The probability spin prior changes the original free energy only by
the deterministic \(n\log2\). The completion is an independent
centered Gaussian of variance \(T/2\), whose scaled \(L^2\) norm is
\(O(n^{-1/6})\) and whose scaled variance is \(T/(2n^{1/3})\).
It follows that the uncompleted scaled variance has the same positive
limit \(v\). Removing that Gaussian and dividing by the exact
finite-size standard deviation is now justified by Slutsky's theorem.
\end{proof}

\subsection{Identification by finite random-size experiments}

\begin{proposition}[Validity of the law specification]\label{red:prop-identification}
Under Proposition~\ref{prop:increment}, the constants in
Proposition~\ref{def:law} are well defined by an absolutely convergent
series after the indicated expectations. They are precisely the
limiting centered moments of \(Y_n\), with \(a_2>0\).
The standardized moments \(a_k/a_2^{k/2}\) determine one Borel
probability measure uniquely. Each \(a_k\) is also the expectation of
a single integrable random variable constructed from fixed priors and
finitely many Gaussian systems almost surely.
\end{proposition}

\begin{proof}
Fix \(k\), let \(m_k(l)=\mathbb EY_l^k\), and denote its power-rate
limit by \(m_k\). Recall
\(C=\sum_{l\ge2}l^{-2}\),
\(\mathbb P(J=l)=l^{-2}/C\), and
\[
 P_k=\prod_{i=1}^k J^{-1/6}(F^{c,0}-F^{c,i}),
 \qquad t=\log J,
\]
where the completed systems are conditionally independent given
\(J\). Conditional on \(J=l\) and copy zero, the other \(k\)
copies average the product to the \(k\)-th centered power of copy
zero. Thus
\[
 \mathbb E[P_k\mid J=l]=m_k(l).
\]
The uniform moment bounds proved above also control every fixed
absolute moment of \(P_k\) conditional on \(J=l\), uniformly in \(l\).

For \(\operatorname{Re}z>0\), put
\[
 B_k(z)=Cz\,\mathbb E[P_k e^{(1-z)t}]
       =z\sum_{l\ge2}m_k(l)l^{-1-z}.
\]
If \(|m_k(l)-m_k|\le C_k l^{-a}\) for some \(a>0\), the difference
series is holomorphic on \(\operatorname{Re}z>-a\), locally
absolutely convergent there. The constant part also extends across
zero. Indeed, for \(\operatorname{Re}z>0\), the difference between
the sum and its integral is
\[
 \sum_{l\ge2}\left(l^{-1-z}-\int_l^{l+1}r^{-1-z}\,dr\right)
       -\int_1^2 r^{-1-z}\,dr.
\]
On a compact subset of \(\operatorname{Re}z>-1\), each summand is
bounded by a constant times \(l^{-2-\operatorname{Re}z}\).
The discrepancy therefore converges locally uniformly and is
holomorphic there. Since
\(z\int_1^\infty r^{-1-z}\,dr=1\), multiplication by \(z\)
removes the integral's pole and gives value one at zero. Combining
this with the difference series, \(B_k\) extends to
\(\operatorname{Re}z>-c_k\) for some \(c_k>0\), with
\(B_k(0)=m_k\).

Its Taylor series about \(z=1\) has radius strictly greater than one.
To evaluate it at zero write \(z=1-w\). Near \(w=0\), differentiation
under the expectation is legitimate: for every fixed \(h\) and
\(0<\delta<1\), the conditional absolute-moment bound on \(P_k\)
and \(\sum_{l\ge2}l^{-2+\delta}(\log l)^h<\infty\) give a
dominating integrable function on \(|w|\le\delta\).
The coefficient of \(w^h\)
in \((1-w)e^{wt}\) is
\[
 \frac{t^h}{h!}
 -{\bf1}_{\{h\ge1\}}\frac{t^{h-1}}{(h-1)!},
\]
with the second term omitted at \(h=0\). Hence
\[
 B_k(0)=C\sum_{h=0}^\infty
 \mathbb E\!\left[P_k\left(
  \frac{t^h}{h!}
  -{\bf1}_{\{h\ge1\}}\frac{t^{h-1}}{(h-1)!}
 \right)\right]=a_k,
\]
and the scalar series is absolutely convergent. The expectation is
taken before the infinite sum. Indeed, at each fixed \(J\) the
bracket series equals \(e^t-e^t=0\); averaging that pointwise sum
would therefore fail to recover the positive second moment.
The limiting law from Proposition~\ref{red:prop-moments} proves
existence of the measure with these moments, its positive variance
proves \(a_2>0\), and its exponential moment bound proves uniqueness.

The final construction uses randomized single-term unbiased estimation
with replication conditional on the sampled level. For this method, see
\cite[Section~2, equation~(10), and Section~4]{RheeGlynn2015}.
Here the levels index Taylor coefficients; the integrability and
almost-sure finiteness proved below do not assert finite expected work
or import the SDE complexity guarantees of that reference.

For the final assertion, draw an independent integer \(H\ge0\) with
\(\mathbb P(H=h)=2^{-h-1}\). Conditional on \(H=h\), repeat the whole
random-size experiment independently \(256^h\) times, and let
\(\overline A_{k,h}\) be the sample average of
\[
 A_{k,h}=CP_k\left(
   \frac{t^h}{h!}
   -{\bf1}_{\{h\ge1\}}\frac{t^{h-1}}{(h-1)!}\right).
\]
Return \(W_k=2^{H+1}\overline A_{k,H}\).
The estimate \(t^h/h!\le4^h e^{t/4}\), the prior proportional to
\(l^{-2}\), and the uniform conditional moments of \(P_k\) give
\(\|A_{k,h}\|_2\le C_k4^h\), including the shifted term.
Independence of the repetitions yields
\[
 \mathbb E\left|
  \overline A_{k,h}-\mathbb EA_{k,h}\right|
 \le\frac{C_k4^h}{\sqrt{256^h}}=C_k4^{-h}.
\]
The means \(\mathbb EA_{k,h}\) are absolutely summable by the Taylor
argument. Consequently
\[
 \mathbb E|W_k|
 =\sum_{h\ge0}\mathbb E|\overline A_{k,h}|<\infty,
 \qquad
 \mathbb EW_k=\sum_{h\ge0}\mathbb EA_{k,h}=a_k.
\]
All priors are fixed and every experiment uses finitely many finite-size
Gaussian systems almost surely. Neither a limiting random input nor a
choice of subsequence enters the specification.
\end{proof}
